\documentclass[11pt,letterpaper]{article}
\usepackage[T1]{fontenc}
\usepackage{lmodern,amsmath,amssymb,amsthm,microtype,booktabs}
\usepackage[margin=1in]{geometry}
\usepackage[hidelinks]{hyperref}
\newtheorem{theorem}{Theorem}
\newtheorem{proposition}{Proposition}
\newcommand{\Z}{\mathbb Z}
\newcommand{\Q}{\mathbb Q}
\newcommand{\F}{\mathbb F}
\title{Prime-seven direction catalogues for cyclotomic unit-distance graphs}
\author{Aafi Khan Malek\\ \small Independent researcher, Ahmedabad, India\\
\small\href{mailto:aafimalek2032@gmail.com}{aafimalek2032@gmail.com}}
\date{28 November 2026}
\begin{document}
\maketitle
\begin{abstract}
Let $\zeta$ be a primitive $42$nd root of unity. Unit-distance
graphs whose edge displacements lie in $\Z_{(7)}[\zeta]$ are
four-colorable. In $\Z[\zeta,1/7]$,
we classify the norm-one units using a single first-order pole
direction from $G_{449}$. There are exactly $273$ unoriented unit
directions with coefficient denominator dividing seven. Exact
enumeration using this catalogue verifies that adjoining any
single point of $\tfrac17\Z[\zeta]$ to the published $G_{449}$
vertex set still gives a five-colorable graph. These restrictions
do not change the chromatic bounds for the plane.
\end{abstract}

\section{Four colors away from the prime seven}

\begin{theorem}
If every edge displacement of a unit-distance graph lies in
$\Z_{(7)}[\zeta]$, then the graph is four-colorable.
\end{theorem}

\begin{proof}
Since $\Phi_{42}(X)\Phi_6(X)=\Phi_6(X^7)$, reduction modulo seven
identifies $\Phi_{42}$ with $(X^2-X+1)^6$. Its two roots in $\F_7$
are $3$ and $5$, which are inverses. They give evaluations
$\rho_3,\rho_5:\Z_{(7)}[\zeta]\to\F_7$: the relation
$\Phi_{42}=0$ holds at each root, and all denominators are
invertible. The two evaluations are interchanged by complex
conjugation. For a unit displacement $w$, the identity
$w\overline w=1$ gives
$\rho_3(w)\rho_5(w)=1$. Its trace $\rho_3(w)+\rho_5(w)$ therefore
belongs to $\{1,2,5,6\}$, by testing the six nonzero residues.

Group the value of $\rho_3(z)+\rho_5(z)$ into
$\{0,3\}$, $\{1,4\}$, $\{2,5\}$, and $\{6\}$.
The only differences within a class are $0$ and $\pm3$, none of
which is an allowed edge trace. For points outside the ring,
translate each additive coset independently; there are no edges
between different cosets.
\end{proof}

\paragraph{The residue graph.}
The finite graph on $\F_7^2$ whose edge steps $(a,b)$ satisfy
$ab=1$ needs four colors. With basis $u=(3,5),v=(5,3)$ the six
steps are $\pm u,\pm v,\pm(u+v)$. In any three-coloring the
triangles on $x,x+u,x+u+v$ and $x,x+v,x+u+v$ force the same
color at $x+u$ and $x+v$. Thus a color depends only on the
sum of the two basis coordinates. The colors at every three
consecutive sums would be distinct, making the sequence
periodic with period three. This is impossible on the cycle
of length seven. Sharpness here is for the residue graph only.

\newpage
\section{Directions with denominators supported at seven}

Put $K=\Q(\zeta)$, with ring of integers $\mathcal O_K=\Z[\zeta]$.
The two primes $P_3=(7,\zeta-3)$ and $P_5=(7,\zeta-5)$ over
seven have ramification index 6, by the factorization of
$\Phi_{42}$ modulo seven. Normalize their valuations to have
value group $\Z$. The exact coefficient audit of the published
$G_{449}$ coordinates finds an edge direction
$u=p(\zeta)/7$, with $p$ of degree below 12, whose numerator
coefficients in increasing order are
\[
(-8,-6,5,5,5,-3,-2,4,7,1,-2,-4).
\]
The norm numerator is exactly $49$ modulo $\Phi_{42}$, and
$p\bmod 7$ has root multiplicities $(5,6)$ at $(3,5)$.

For the valuation at $P_3$, the element $\zeta-3$ is a local
uniformizer: together with $7$ it generates the prime, while
$v_{P_3}(7)=6$. In the expansion of $p$ in powers of $\zeta-3$,
the terms of degree below five have coefficients divisible by
seven, and the degree-five coefficient is not. The unique
least valuation is five. Hence $v_{P_3}(u)=-1$, and by conjugation
and $u\overline u=1$ we have $v_{P_5}(u)=1$.

\begin{theorem}
Every $w\in\mathcal O_K[1/7]$ with $w\overline w=1$ has the
unique expression $w=\zeta^k u^n$ for
$k\in\Z/42\Z$ and $n\in\Z$. Its least common coefficient
denominator is $7^{\lceil |n|/6\rceil}$.
\end{theorem}

\begin{proof}
The valuations of $w$ at $P_3$ and $P_5$ are opposites. Take
$n=-v_{P_3}(w)$. Then the quotient $w/u^n$ has zero valuation
at both primes over seven and is integral at every other
prime, so it is an algebraic integer. It still has conjugate
norm one. All of its Galois conjugates have modulus one,
because complex conjugation commutes with every automorphism
of $K$. Kronecker's theorem makes it a root of unity. Every
root of unity in $K$ has order dividing 42: otherwise that
root together with $\zeta$ would generate a cyclotomic field
of order a proper multiple of 42 and strictly larger degree.
This proves the expression, and valuations give uniqueness.

Since $\mathcal O_K$ has the integral power basis, multiplying
by $7^h$ clears the coefficient denominator exactly when it
makes every valuation nonnegative. Only the two primes over
seven matter, and both receive $6h$. Thus the least $h$ is
$\lceil |n|/6\rceil$.
\end{proof}

The norm-one catalogue with $|n|\le B$ has $21(2B+1)$
unoriented directions. In particular there are 63 directions
for poles of order at most one and 273 for denominators
dividing seven.

\newpage
\section{Exact checks on the published vertex set}

We used the pinned coordinate and sparse6 files of Heule and
Przybocki~\cite{hp}. The declared denominator is 49, but the
integer coefficient numerators share a factor of seven. The
eight cosets modulo $\Z[\zeta]$ and their internal edge and
triangle counts are as follows, with the two equal-size
cosets ordered by their least input vertex index.
\[
\begin{array}{lrrrrrrrr}
\toprule
\text{Size} &1&49&49&56&63&70&77&84\\
\text{Internal edges}&0&105&105&126&147&168&203&238\\
\text{Internal triangles}&0&30&9&27&27&23&40&39\\
\bottomrule
\end{array}
\]
The supplied audit gives an explicit three-coloring of every
non-singleton layer; the triangles make each layer exactly
three-chromatic. A related general result for directions
whose angles differ by rational multiples of $\pi$ is given
by Davidson~\cite{davidson}. Exactly 1,092 edges are integral
displacements and 1,106 connect distinct cosets. All 49
observed unoriented directions occur in the first-order pole
catalogue. Applying the full denominator-seven catalogue to
the existing points recovers exactly all 2,198 published
edges and no others.

\begin{proposition}
The induced unit-distance graph on the $449$ published points
together with any one point in $\tfrac17\Z[\zeta]$ is
five-colorable.
\end{proposition}

\begin{proof}
Translating the 273 directions with both orientations from
the 449 points produces 229,320 distinct outside neighbors;
all other outside grid points have no edge to the original
graph. Under the published five-coloring, exactly 34 of these
outside neighbors see all five colors. Every other point can
be given a missing color. For the 34 exceptions the companion
certificate supplies a valid five-coloring of each one-point
extension, and additionally a five-coloring on all 483 points at
once with 2,395 exact edges. The certificate checker recomputes
the candidate points and all relevant edges by integer arithmetic.
\end{proof}

\paragraph{Reproducibility and scope.}
The repository supplies the catalogue generator, both direction
tables, and the coloring witnesses. The Init-only Lean 4.36.0
package checks the finite residue graph, the factorization
modulo seven, and the four-color rule for an arbitrary graph
given exact polynomial norm certificates. Deriving these
certificates from arbitrary complex unit edges, the valuation
classification, and the G449 data audit are not formalized in
Lean here. The results do not exclude larger extensions or improve
the plane bounds.

\paragraph{Contributions.}
Aafi Khan Malek selected and directed the research. ChatGPT assisted
with proof development, software, formalization and writing. Codex
was not used. This note has not been peer reviewed.

\begin{thebibliography}{9}
\bibitem{hp} M. J. H. Heule and B. Przybocki,
\emph{From two heptagrams to a 449-vertex 5-chromatic unit-distance graph},
\href{https://arxiv.org/abs/2611.28518}{arXiv:2611.28518}, 2026.
\bibitem{davidson} J. Davidson,
\emph{Coloring monomial Cayley graphs: the one-parameter case},
\href{https://arxiv.org/abs/2610.09474}{arXiv:2610.09474}, 2026.
\end{thebibliography}
\end{document}
